\documentclass[11pt]{article}
\usepackage[margin=2.5cm]{geometry}
\usepackage{amsmath,amssymb,bm}
\usepackage{booktabs,longtable,array}
\usepackage[hidelinks]{hyperref}

\newcommand{\dd}{\mathrm{d}}
\newcommand{\ii}{\mathrm{i}}
\newcommand{\T}{\mathsf{T}}
\newcommand{\Ifft}{\mathcal{F}^{-1}}
\newcommand{\Real}{\operatorname{Re}}
\newcommand{\Imag}{\operatorname{Im}}
\newcommand{\clip}{\operatorname{clip}}
\newcommand{\vsv}{V_{SV}}
\newcommand{\vsh}{V_{SH}}
\newcommand{\vpv}{V_{PV}}
\newcommand{\vph}{V_{PH}}

\title{Step 2: Joint Love--Rayleigh multimode array waveform inversion\\
\large\texttt{RL\_array\_multimode\_waveform\_inv} (drivers \texttt{D2}--\texttt{D4})}
\date{}

\begin{document}
\maketitle

\noindent This note describes the second step of the two-step inversion:
a transdimensional Bayesian (reversible-jump Markov chain Monte Carlo,
rjMcMC) inversion of multimode Love and Rayleigh waveforms recorded across
an ocean-bottom array for the depth profiles of the vertically polarized
shear velocity $\vsv(z)$ and the radial anisotropy $\xi(z)$. The reference
model is the isotropic model produced by step 1 (an inversion of
fundamental-mode Rayleigh phase velocities). Every symbol is defined where
it first appears and is collected in the table of
Section~\ref{sec:notation}. Default values are those of
\texttt{parameters\_RL.m} and \texttt{functions/sub\_setup\_records.m}.

\tableofcontents

%======================================================================
\section{Overview}
%======================================================================

The unknowns are split into three groups:
\begin{enumerate}
\item \textbf{Structure} $\mathbf{m}$: the profiles $\vsv(z)$ and
  $\xi(z)=(\vsh/\vsv)^2$ between the Moho and $z_{\max}=350$~km, each a
  B-spline curve with a variable number of control points.
  $\vsv$ is the vertically polarized and $\vsh$ the horizontally polarized
  shear-wave speed, and $z$ is depth below the seafloor.
\item \textbf{Noise hyperparameters} $\bm{\sigma}=(\sigma_q)$ and
  $\mathbf{r}=(r_q)$: a noise standard deviation and a lag-one correlation
  for each of the four window types $q$ (Section~\ref{sec:windows}).
\item \textbf{Nuisance parameters} $\mathbf{x}_r$ of every record $r$: for
  each modeled branch, a smooth complex spectral correction $B_n(\omega)$,
  a frequency-dependent time shift $\tau_n(\omega)$ and a
  frequency-dependent arrival-angle anomaly $\delta\theta_n(\omega)$
  (Section~\ref{sec:forward}).
\end{enumerate}
The structure and the noise hyperparameters are sampled by Markov chain
Monte Carlo. The nuisance parameters are integrated out of the likelihood
(Section~\ref{sec:like}): $B_n$ exactly, because it enters linearly, and
$\tau_n$ and $\delta\theta_n$ by linearizing the waveforms about a
\emph{linearization point}. Each chain state carries its own linearization
point, which is moved to the conditional posterior mean every
$n_{\rm relin}$ iterations (Section~\ref{sec:relin}). The sampler is either
a set of independent chains (``plain'') or parallel tempering
(Section~\ref{sec:mcmc}).

%======================================================================
\section{Data}\label{sec:data}
%======================================================================

\subsection{Geometry}
The events are indexed $e=1,\dots,N_e$ ($N_e=15$ real 2012 Pacific-rim
earthquakes with Global CMT moment tensors, $M_w$ 6.4--7.8, 33--86$^\circ$
from the array). The stations are indexed $j=1,\dots,N_s$ ($N_s=16$
NoMelt ocean-bottom seismometers). Station latitudes and longitudes
$(\phi_j,\lambda_j)$ are projected to local Cartesian coordinates (km,
east and north) about the array center $(\phi_0,\lambda_0)$, the mean
station position:
\begin{equation}
  x_j = (\lambda_j-\lambda_0)\,a_\circ\cos\phi_0,\qquad
  y_j = (\phi_j-\phi_0)\,a_\circ,\qquad a_\circ = 111.195~\text{km/deg}.
\end{equation}
$\Delta_e$ is the great-circle distance (km) from event $e$ to the array
center, and $\Delta_{ej}$ the distance to station $j$. The propagation
azimuth $\varphi_e$ (degrees clockwise from north) is the back azimuth from
the array center to the event plus $180^\circ$.

The wave type is $w\in\{L,R\}$. Love waves ($L$) are recorded on the
transverse component and Rayleigh waves ($R$) on the vertical component. A
\emph{record} $r=(e,w)$ is one event on one component across all stations,
so there are $N_r=2N_e=30$ records.

\subsection{Synthetic data}
The ``observed'' seismograms $u^{\rm obs}_{ej}(t)$ are \texttt{idagrn6}
normal-mode synthetics for the true model (PA5) with the mode branches
$n=0,\dots,N_b^{\rm data}-1$, where $N_b^{\rm data}=16$, at sampling
interval $\Delta t=4$~s. The inversion models only the branches
$n=0,\dots,n_m-1$, with $n_m=8$. Branches 8--15 are present in the data
but not in the model.

\subsection{Time segment, taper and band}\label{sec:proc}
Each record is cut to a segment of $N=800$ samples, $t_i=t^0_r+i\,\Delta t$
for $i=0,\dots,N-1$, starting at
\begin{equation}
 t^0_r=\Delta t\;\mathrm{round}\!\left[\frac{\Delta_e/U^{\rm ov}_w-150~\text{s}}{\Delta t}\right],
\end{equation}
where $U^{\rm ov}_w$ is the overtone-window group-velocity bound of
Section~\ref{sec:windows}. The segment is multiplied by a Tukey window
$h(t_i)$ with a 10\% taper and Fourier transformed. The discrete
frequencies are $f_k=k/(N\Delta t)$ and the angular frequencies
$\omega_k=2\pi f_k$. Only the band
\begin{equation}
 \mathcal{B}=\{k: f_1<f_k<f_4\},\qquad (f_1,f_2,f_3,f_4)=(1/112,\,1/100,\,1/20,\,1/17)~\text{Hz}
\end{equation}
is used, with $N_f=|\mathcal{B}|$ frequencies. A cosine band-pass taper
$H(f)$ is applied: $H=1$ for $f_2\le f\le f_3$; on $(f_1,f_2)$ it rises as
$H=\tfrac12-\tfrac12\cos[\pi(f-f_1)/(f_2-f_1)]$; and on $(f_3,f_4)$ it
falls as $H=\tfrac12+\tfrac12\cos[\pi(f-f_3)/(f_4-f_3)]$.

The spectrum of a trace $u$ in the band is
\begin{equation}
  \hat u(\omega_k)=H(f_k)\sum_{i=0}^{N-1}h(t_i)\,u(t_i)\,e^{-2\pi\ii\,ki/N},\qquad k\in\mathcal{B}.
\end{equation}
The inverse operator $\Ifft$ takes a one-sided band spectrum $X(\omega_k)$,
$k\in\mathcal{B}$, to a real time series:
\begin{equation}\label{eq:ifft}
  \Ifft[X](t_i)=2\,\Real\Big\{\frac1N\sum_{k\in\mathcal{B}}X(\omega_k)\,e^{2\pi\ii\,ki/N}\Big\}.
\end{equation}
The processed noise-free data trace is $u^{\rm clean}_{ej}=\Ifft[\hat u^{\rm obs}_{ej}]$.

\subsection{Windows from reference group velocities}\label{sec:windows}
$U_{n}^{w,\rm ref}(T)$ is the MINEOS group velocity of branch $n$ in the
reference (step-1) model at period $T$, evaluated at the dispersion nodes
(Section~\ref{sec:disp}). For each wave type, three bounds (km/s) are
defined:
\begin{align}
  U^{\rm b}_w &= \max_T U^{w,\rm ref}_0(T)+0.2, \\
  U^{\rm ov}_w &= \max_{T,\;1\le n<n_m} U^{w,\rm ref}_n(T)+0.6\quad\text{(or \texttt{opt.U\_start})}, \\
  U^{\rm f}_w &= \min_T U^{w,\rm ref}_0(T)-0.35 .
\end{align}
For the step-1 model, $U^{\rm ov}=8.30$ (Love) and $9.24$~km/s (Rayleigh).
With \texttt{opt.U\_start} a slower start (e.g.\ 6~km/s) leaves out the
early, body-wave-like energy (S, SS, \dots), which in mode terms is
carried mostly by the unmodeled branches $n\ge n_m$.

\emph{Single window} (\texttt{opt.windows = 'single'}, default): at station
$j$ of record $r$ the data are one surface-wave window
$\Delta_{ej}/U^{\rm ov}_w\le t<\Delta_{ej}/U^{\rm f}_w+60$~s, containing the
overtones and the fundamental mode. Window types are labelled by wave
type, $q\in\{1,3\}=\{$Love, Rayleigh$\}$.

\emph{Split windows} (\texttt{opt.windows = 'split'}): two windows,
\begin{itemize}
\item the \emph{overtone window} $\Delta_{ej}/U^{\rm ov}_w\le t<\Delta_{ej}/U^{\rm b}_w$;
\item the \emph{fundamental-mode window} $\Delta_{ej}/U^{\rm b}_w\le t<\Delta_{ej}/U^{\rm f}_w+60$~s,
\end{itemize}
labelled $q\in\{1,2,3,4\}=\{$Love overtone, Love fundamental, Rayleigh
overtone, Rayleigh fundamental$\}$.

The samples are the same in both cases; the choice only sets how many
noise classes there are (Section~\ref{sec:like}). Separate classes let the
overtone part, which also contains the unmodeled branches, have its own
noise level, so it does not set the weight of the better-fit fundamental
mode, nor the reverse. $q(i)$ is the window type of sample $i$. A
\emph{segment} is a run of consecutive samples of one station and one
window type. In the single-window case the samplers keep the four noise
slots and tie $(\sigma_1,r_1)=(\sigma_2,r_2)$ and $(\sigma_3,r_3)=(\sigma_4,r_4)$.

\subsection{Data vector, normalization and noise}
The data vector $\mathbf{d}_r\in\mathbb{R}^{N_{d,r}}$ of record $r$ stacks
the windowed samples of all stations. It is normalized by the peak
amplitude of the noise-free windowed record,
$s_r=\max_i|u^{\rm clean}_{r,i}|$, as is usual with real data:
\begin{equation}
  \mathbf{d}_r=\mathbf{u}^{\rm clean}_r/s_r+\bm{\epsilon}_r .
\end{equation}
The noise $\bm{\epsilon}_r$ is white Gaussian noise passed through the same
band-pass $H$, scaled to a standard deviation $\sigma^{\rm true}=0.03$
(3\% of the record peak). Because it is band-limited, the noise is
correlated in time; its empirical lag-one correlation is
$\hat r_q\approx0.67$. $N_d=\sum_rN_{d,r}=181{,}172$ samples.

\subsection{Reference spectra}
$U^{\rm ref}_{rjn}(\omega_k)$ is the band spectrum (same segment, Tukey
window, taper $H$ and normalization $s_r$) of the \texttt{idagrn6\_mask}
single-branch synthetic of branch $n$ in the \emph{reference} model at
station $j$. It contains the source excitation, the propagation from the
source to station $j$ in the reference model and the instrument-free
displacement response. For plotting, the reference prediction is
$\Ifft[\sum_n U^{\rm ref}_{rjn}]$, and the PREM prediction is the processed
full PREM synthetic (all 16 branches) divided by $s_r$.

%======================================================================
\section{Model parameterization}\label{sec:model}
%======================================================================

\subsection{Transdimensional B-spline profiles}
$z$ is depth below the seafloor (the water depth is $z_w$). The inverted
depth range is $z_{\min}=7$~km (the Moho) to $z_{\max}=350$~km. Each
profile $p\in\{\vsv,\xi\}$, and $p=\eta$ when $\eta$ is a model parameter
(Section~\ref{sec:derived}), is described by:
\begin{itemize}
\item $k_p$ interior control points $(z^p_i,v^p_i)$, $i=1,\dots,k_p$, with $z_{\min}<z^p_i<z_{\max}$;
\item a free top end value $v^p_{\rm top}$ at $z_{\min}$;
\item a bottom end value $v^p_{\rm bot}$ at $z_{\max}$, \emph{fixed} to the
  reference model there (e.g.\ $\xi_{\rm bot}=1$).
\end{itemize}
Sorted by depth, the $n_p=k_p+2$ control points $\mathbf{P}_l=(Z_l,V_l)$,
$l=1,\dots,n_p$, define a clamped parametric B-spline curve of degree
$\varrho=\min(3,n_p-1)$:
\begin{equation}
  \big(z(s),\,v(s)\big)=\sum_{l=1}^{n_p}N_{l,\varrho}(s)\,(Z_l,V_l),\qquad s\in[0,1].
\end{equation}
Here $N_{l,\varrho}(s)$ are the Cox--de Boor basis functions on the clamped
uniform knot vector
\begin{equation}
 \Big(\underbrace{0,\dots,0}_{\varrho+1},\;\tfrac{1}{n_p-\varrho},\;\tfrac{2}{n_p-\varrho},\dots,\tfrac{n_p-\varrho-1}{n_p-\varrho},\;\underbrace{1,\dots,1}_{\varrho+1}\Big).
\end{equation} Because the $Z_l$ are sorted, $z(s)$ is monotonic and
the curve is a single-valued, $C^2$ profile $v(z)$, evaluated by
interpolation. Clamping makes the profile pass through the end values.
Adding or removing a control point changes the curve only locally.

The model is
$\mathbf{m}=\{k_v,\{z^v_i,v^v_i\},v^v_{\rm top};\;k_x,\{z^x_i,v^x_i\},v^x_{\rm top}\}$,
plus $\{k_e,\{z^e_i,v^e_i\},v^e_{\rm top}\}$ for $\eta$ when it is free.
It gives $\vsv(z)$, $\xi(z)$ (and $\eta(z)$) on $[z_{\min},z_{\max}]$. Outside this range
(water, crust, $z>350$~km) the model is the reference model.

\subsection{Derived parameters}\label{sec:derived}
\begin{equation}
  \vsh(z)=\vsv(z)\sqrt{\xi(z)} .
\end{equation}
$\vpv$ and $\vph$ are the vertically and horizontally propagating P-wave
speeds and $\rho$ is density. They follow the shear velocities with the
ratios of the reference model (superscript 0):
\begin{equation}\label{eq:scaling}
  \vpv=R_{pv}\,\vsv,\quad \vph=R_{ph}\,\vsh,\quad \rho=R_\rho\,\vsv,\qquad
  R_{pv}=\frac{\vpv^0}{\vsv^0},\;R_{ph}=\frac{\vph^0}{\vsh^0},\;R_\rho=\frac{\rho^0}{\vsv^0}.
\end{equation}
The anisotropic parameter $\eta$ (the fifth Love parameter,
$\eta=F/(A-2L)$) is held at the reference model by default,
$\eta=\eta^0$ (\texttt{opt.eta\_free = false}). With
\texttt{opt.eta\_free = true} it is a third transdimensional B-spline
profile $\eta(z)$ (Section~\ref{sec:prior}).

\subsection{Dispersion}\label{sec:disp}
$c^w_n(T;\mathbf{m})$ is the phase velocity of branch $n$ of wave type $w$
at period $T$. It is evaluated at $N_T=10$ dispersion nodes $T_\ell$,
log-spaced from $0.95\times17$ to $1.02\times112$~s, by linearizing about
the reference model with the MINEOS kernels of every parameter:
\begin{equation}\label{eq:disp}
  c^w_n(T_\ell;\mathbf{m})=c^{w,\rm ref}_n(T_\ell)+\sum_{i\in\mathcal{N}}\;
  \sum_{X\in\mathcal{X}}G^{w,X}_{n,\ell i}\,\delta X(z_i),\qquad
  \mathcal{X}=\{\vsv,\vsh,\vpv,\vph,\rho,\eta\}.
\end{equation}
The terms are:
\begin{itemize}
\item $\mathcal{N}$: the reference card nodes with
  $z_{\min}\le z_i\le z_{\max}$ (only the mantle-side node at the Moho);
\item $c^{w,\rm ref}_n$: the MINEOS phase velocity of the reference model;
\item $\delta X(z_i)=X(z_i;\mathbf{m})-X^0(z_i)$: the perturbation of
  parameter $X$ from the reference model at node $i$, in km/s for the
  velocities, g/cm$^3$ for $\rho$ and dimensionless for $\eta$. The
  perturbations follow from the model through Eq.~\eqref{eq:scaling}:
  \begin{gather*}
   \delta\vsv=\vsv(\mathbf{m})-\vsv^0,\qquad \delta\vsh=\vsv(\mathbf{m})\sqrt{\xi(\mathbf{m})}-\vsh^0,\\
   \delta\vpv=R_{pv}\vsv(\mathbf{m})-\vpv^0,\qquad \delta\vph=R_{ph}\vsh(\mathbf{m})-\vph^0,\qquad
   \delta\rho=R_\rho\vsv(\mathbf{m})-\rho^0,
  \end{gather*}
  and $\delta\eta=\eta(\mathbf{m})-\eta^0$, which is zero unless $\eta$ is
  free;
\item $G^{w,X}_{n,\ell i}$: the kernel of parameter $X$,
  \begin{equation}
   G^{w,X}_{n,\ell i}=K^{w,n}_{X}(T_\ell,r_i)\,\Delta r_i\cdot10^3\quad(X\ne\eta),\qquad
   G^{w,\eta}_{n,\ell i}=K^{w,n}_{\eta}(T_\ell,r_i)\,\Delta r_i .
  \end{equation}
\end{itemize}
$K^{w,n}_{X}(T,r)$ is the MINEOS Fréchet kernel of phase velocity with
respect to parameter $X$ at radius $r$, in the units where
$\delta c\,[\text{km/s}]=\sum_iK_X(r_i)\,\Delta r_i[\text{m}]\,\delta X$, with
$\delta X$ in m/s, kg/m$^3$ or dimensionless ($\eta$). $\Delta r_i$ is the
node weight (the centered difference of node radii, in m). The factor
$10^3$ converts to km/s and g/cm$^3$. Because
$\delta\vpv=R_{pv}\delta\vsv$, $\delta\vph=R_{ph}\delta\vsh$ and
$\delta\rho=R_\rho\delta\vsv$, Eq.~\eqref{eq:disp} is identical to
grouping the Vp and density kernels with the shear kernels (verified to
$10^{-14}$~km/s).

Love waves (toroidal modes) depend on $\vsh$, $\vsv$ and $\rho$; their
$\vpv$, $\vph$ and $\eta$ kernels are zero. Rayleigh waves depend on all
six parameters, so both wave types constrain $\xi$, and only Rayleigh waves
constrain $\eta$. The phase velocities are interpolated from the nodes
$T_\ell$ to the band frequencies with a monotone cubic (PCHIP) interpolant
in $\ln T$. A frequency where branch $n$ has no node on either side is
masked: $\chi^w_n(\omega)=0$ there and 1 elsewhere. A model is rejected
if any phase velocity needed at an existing branch is not finite.

%======================================================================
\section{Forward model of the waveforms}\label{sec:forward}
%======================================================================

\subsection{Modal spectra}
For record $r=(e,w)$, station $j$ and modeled branch $n$, the predicted
band spectrum is
\begin{equation}\label{eq:forward}
  U_{rjn}(\omega)=B_{rn}(\omega)\,U^{\rm ref}_{rjn}(\omega)\,
  e^{-\ii\omega\tau_{rn}(\omega)}\,
  \exp\!\left\{-\ii\omega\left[\frac{\Delta D_j\big(\varphi_e+\delta\theta_{rn}(\omega)\big)}{c^w_n(\omega;\mathbf{m})}
  -\frac{\Delta D_j(\varphi_e)}{c^{w,\rm ref}_n(\omega)}\right]\right\}\chi^w_n(\omega).
\end{equation}
The plane-wave path difference of station $j$ relative to the array center
for a wave propagating toward azimuth $\varphi$ is
\begin{equation}
  \Delta D_j(\varphi)=x_j\sin\varphi+y_j\cos\varphi .
\end{equation}
The factors in Eq.~\eqref{eq:forward} mean:
\begin{itemize}
\item $U^{\rm ref}_{rjn}$ already contains the full reference-model
  propagation to station $j$.
\item The last exponential replaces the reference phase velocity across
  the array by the model phase velocity. This term carries the structural
  information.
\item $\tau_{rn}(\omega)$ is the time shift of branch $n$ common to all
  stations. It absorbs the difference between the true and reference path
  from the source to the array center, which is not part of the 1-D array
  model.
\item $\delta\theta_{rn}(\omega)$ is the arrival-angle anomaly in degrees,
  common to all stations: off-great-circle propagation.
\item $B_{rn}(\omega)$ is a smooth complex correction of amplitude and phase:
  source excitation, focusing and attenuation errors.
\end{itemize}

The predicted data are the sum over the modeled branches, transformed with
Eq.~\eqref{eq:ifft} and sampled in the windows:
\begin{equation}
  \mathbf{u}_r(\mathbf{m},\mathbf{x}_r)=\Big[\Ifft\Big[\sum_{n=0}^{n_m-1}U_{rjn}\Big](t_i)\Big]_{(j,i)\in\text{windows}} .
\end{equation}

\subsection{Frequency bases of the nuisance functions}
Let $\zeta(\omega)=2\,\dfrac{\ln\omega-\ln\omega_{\min}}{\ln\omega_{\max}-\ln\omega_{\min}}-1\in[-1,1]$
be normalized log frequency over the band, and $P_a$ the Legendre
polynomials: $P_0=1$, $P_1=\zeta$, $P_2=\tfrac32\zeta^2-\tfrac12$,
$P_3=\tfrac52\zeta^3-\tfrac32\zeta$. Then
\begin{align}
  B_{rn}(\omega)&=\sum_{k=1}^{K}\beta_{rnk}\,\Phi_k(\omega),\qquad \beta_{rnk}\in\mathbb{C},\\
  \tau_{rn}(\omega)&=\sum_{a=1}^{K_\tau}t_{rna}\,P_{a-1}\big(\zeta(\omega)\big)\quad\text{(s)},\\
  \delta\theta_{rn}(\omega)&=\sum_{a=1}^{K_\theta}\theta_{rna}\,P_{a-1}\big(\zeta(\omega)\big)\quad\text{(deg)} .
\end{align}
$\Phi_k$ are $K=3$ piecewise-linear ``hat'' functions of $\ln\omega$. Their
peaks are at $K$ equally spaced nodes $\ln\omega_{\min}=\nu_1<\dots<\nu_K=\ln\omega_{\max}$:
$\Phi_k=\max(0,1-|\ln\omega-\nu_k|/(\nu_2-\nu_1))$, and they sum to 1.

The defaults are $K_\tau=3$ (\texttt{opt.spec\_Kt}: constant, linear and
quadratic in $\zeta$) and $K_\theta=2$. A frequency-dependent $\tau$ is
needed for the overtones. Their path delays relative to the reference,
$\Delta_e\,[1/c^{\rm true}_n(\omega)-1/c^{\rm ref}_n(\omega)]$, vary with
frequency because each period samples the unmodeled structure differently.

The nuisance vector of record $r$ is
\begin{equation}
 \mathbf{x}_r=\big[\underbrace{\Real\beta_{rnk},\Imag\beta_{rnk}}_{2n_mK}\;\big|\;\underbrace{t_{rna}}_{n_mK_\tau}\;\big|\;\underbrace{\theta_{rna}}_{n_mK_\theta}\big],
\end{equation}
of length $P=n_m(2K+K_\tau+K_\theta)=88$.

\subsection{Linearized operator}\label{sec:G}
The waveforms are exactly linear in $\beta$ but not in $t$ and $\theta$.
They are linearized about a linearization point
$\bm{\ell}_r=(\mathbf{b}^0_r,\mathbf{t}^0_r,\bm{\theta}^0_r)$, which
holds reference values of $\beta$, $t$ and $\theta$. Define the unscaled
branch spectrum at the linearization point,
\begin{equation}
 S_{rjn}(\omega)=U^{\rm ref}_{rjn}\,e^{-\ii\omega\tau^0_{rn}(\omega)}
 \exp\!\Big\{-\ii\omega\Big[\tfrac{\Delta D_j(\varphi^0_{rn}(\omega))}{c^w_n(\omega;\mathbf{m})}-\tfrac{\Delta D_j(\varphi_e)}{c^{w,\rm ref}_n(\omega)}\Big]\Big\}\chi^w_n(\omega),
\end{equation}
with $\tau^0_{rn}=\sum_at^0_{rna}P_{a-1}$ and
$\varphi^0_{rn}=\varphi_e+\sum_a\theta^0_{rna}P_{a-1}$. The prediction is
then approximated as
\begin{equation}\label{eq:lin}
 \mathbf{u}_r\approx\mathbf{G}_r(\mathbf{m},\bm\ell_r)\,\tilde{\mathbf{x}}_r,\qquad
 \tilde{\mathbf{x}}_r=\big[\Real\bm\beta_r,\Imag\bm\beta_r\;|\;\delta\mathbf{t}_r\;|\;\delta\bm\theta_r\big],
\end{equation}
where $\delta\mathbf{t}_r=\mathbf{t}_r-\mathbf{t}^0_r$ and
$\delta\bm\theta_r=\bm\theta_r-\bm\theta^0_r$. The columns of $\mathbf{G}_r$ are:
\begin{itemize}
\item \emph{Amplitude columns} (exact): with the complex inverse transform
  \[
   a_{nk}(t_i)=\frac1N\sum_{k'\in\mathcal{B}}\Phi_k(\omega_{k'})\,S_{rjn}(\omega_{k'})\,e^{2\pi\ii k'i/N}
  \]
  (sampled in the windows to form $\mathbf{a}_{nk}$), the columns for $\Real\beta_{nk}$ and
  $\Imag\beta_{nk}$ are $2\Real\mathbf{a}_{nk}$ and
  $-2\Imag\mathbf{a}_{nk}$, because
  $\Ifft[\beta\Phi_kS]=\Real\beta\cdot2\Real\mathbf{a}-\Imag\beta\cdot2\Imag\mathbf{a}$.
\item \emph{Time-shift columns}: $\Ifft\big[-\ii\omega P_{a-1}(\zeta)\,\tilde B_{rn}S_{rjn}\big]$.
\item \emph{Arrival-angle columns}:
  $\Ifft\big[-\ii\omega\,c^w_n(\omega)^{-1}\,\partial_\varphi\Delta D_j(\varphi^0_{rn})\,\tfrac{\pi}{180}\,P_{a-1}(\zeta)\,\tilde B_{rn}S_{rjn}\big]$,
  with $\partial_\varphi\Delta D_j=x_j\cos\varphi-y_j\sin\varphi$.
\end{itemize}
In the last two, $\tilde B_{rn}=\sum_k\tilde\beta_{rnk}\Phi_k$ is the
amplitude at which the derivatives are taken. $\tilde\beta_r$ is
re-estimated for the current structure by damped least squares on the
amplitude columns $\mathbf{G}^\beta_r$:
\begin{equation}
 \tilde{\bm\beta}_r=\arg\min_{\bm\beta}\|\mathbf{d}_r-\mathbf{G}^\beta_r\bm\beta\|^2+\lambda\|\bm\beta-\bm\mu_\beta\|^2,\qquad
 \lambda=(\bar\sigma/\sigma_b)^2 .
\end{equation}
Here $\bar\sigma$ is the mean of the four $\sigma_q$, and $\bm\mu_\beta$ and
$\sigma_b$ are the prior mean and standard deviation of
Section~\ref{sec:nuisprior}. Equation~\eqref{eq:lin} is the first-order
(Gauss--Newton) expansion of $\mathbf{u}_r$ about
$(\tilde{\bm\beta}_r,\mathbf{t}^0_r,\bm\theta^0_r)$ in $t$ and $\theta$.

%======================================================================
\section{Likelihood with the nuisance parameters integrated out}\label{sec:like}
%======================================================================

\subsection{Prior of the nuisance parameters}\label{sec:nuisprior}
The prior is Gaussian and independent, $\tilde{\mathbf{x}}_r\sim\mathcal{N}(\bm\mu_r,\mathbf{C}_x)$:
\begin{itemize}
\item $\Real\beta_{rnk}\sim\mathcal{N}(1,\sigma_b^2)$ and
  $\Imag\beta_{rnk}\sim\mathcal{N}(0,\sigma_b^2)$, with $\sigma_b=0.4$
  (no correction a priori).
\item The \emph{total} time-shift coefficients
  $t_{rna}=t^0_{rna}+\delta t_{rna}\sim\mathcal{N}(0,s_{\tau,a}^2)$, with
  $\mathbf{s}_\tau=(15,10,10)$~s (\texttt{opt.tau\_prior},
  \texttt{opt.dtau\_prior}). Hence $\delta t_{rna}$ has mean $-t^0_{rna}$.
\item The total arrival-angle coefficients
  $\theta_{rna}\sim\mathcal{N}(0,s_{\theta,a}^2)$, with
  $\mathbf{s}_\theta=(5,5)^\circ$. Hence $\delta\theta_{rna}$ has mean
  $-\theta^0_{rna}$.
\end{itemize}
Centering the prior on the total value makes it independent of where the
linearization point is.

\subsection{Correlated noise}
The noise of record $r$ is $\mathbf{e}_r\sim\mathcal{N}(\mathbf{0},\mathbf{C}_{d,r})$.
$\mathbf{C}_{d,r}$ is block diagonal over segments, and within a segment of
window type $q$ it is an AR(1) (exponential) covariance:
\begin{equation}
 [\mathbf{C}_{d}]_{il}=\sigma_q^2\,r_q^{|i-l|},\qquad 0\le r_q<1 .
\end{equation}
Its inverse is tridiagonal, so the data and every column of $\mathbf{G}$
are \emph{whitened} segment by segment:
\begin{equation}
 \tilde y_{1}=y_{1},\qquad \tilde y_i=\frac{y_i-r_q\,y_{i-1}}{\sqrt{1-r_q^2}}\quad(i>1\text{ within the segment}).
\end{equation}
The whitened noise is white with variance $\sigma_q^2$, and
\begin{equation}
 \ln|\mathbf{C}_{d,r}|=\sum_q\big[n_{rq}\ln\sigma_q^2+(n_{rq}-m_{rq})\ln(1-r_q^2)\big],
\end{equation}
where $n_{rq}$ is the number of samples and $m_{rq}$ the number of
segments of window type $q$ in record $r$. Below, $\mathbf{d}$ and
$\mathbf{G}$ denote the whitened quantities and $\mathbf{C}_d$ is diagonal
with entries $\sigma_{q(i)}^2$.

\subsection{Marginal likelihood}
With the linear model $\mathbf{d}_r=\mathbf{G}_r\tilde{\mathbf{x}}_r+\mathbf{e}_r$
and the Gaussian prior, the nuisance parameters integrate out exactly:
\begin{equation}
 p(\mathbf{d}_r\mid\mathbf{m},\bm\sigma,\mathbf{r},\bm\ell_r)=\int\mathcal{N}(\mathbf{d}_r;\mathbf{G}_r\tilde{\mathbf{x}},\mathbf{C}_d)\,\mathcal{N}(\tilde{\mathbf{x}};\bm\mu_r,\mathbf{C}_x)\,\dd\tilde{\mathbf{x}}
 =\mathcal{N}\big(\mathbf{d}_r;\mathbf{G}_r\bm\mu_r,\;\mathbf{C}_d+\mathbf{G}_r\mathbf{C}_x\mathbf{G}_r^\T\big).
\end{equation}
It is evaluated with the Woodbury identity, so only a $P\times P$ system
is solved:
\begin{gather}
 \mathbf{A}_r=\mathbf{C}_x^{-1}+\mathbf{G}_r^\T\mathbf{C}_d^{-1}\mathbf{G}_r,\qquad
 \bm\rho_r=\mathbf{d}_r-\mathbf{G}_r\bm\mu_r,\qquad
 \mathbf{g}_r=\mathbf{G}_r^\T\mathbf{C}_d^{-1}\bm\rho_r,\\
 \ln\mathcal{L}_r=-\tfrac12\Big[\bm\rho_r^\T\mathbf{C}_d^{-1}\bm\rho_r-\mathbf{g}_r^\T\mathbf{A}_r^{-1}\mathbf{g}_r
 +\ln|\mathbf{C}_{d,r}|+\ln|\mathbf{C}_x|+\ln|\mathbf{A}_r|+N_{d,r}\ln2\pi\Big].\label{eq:logL}
\end{gather}
Here $\ln|\mathbf{C}_{d,r}|$ includes the correlation term above. The total
log-likelihood is $\ln\mathcal{L}(\mathbf{m},\bm\sigma,\mathbf{r})=\sum_r\ln\mathcal{L}_r$.

The conditional posterior mean of the nuisance parameters is
\begin{equation}\label{eq:xhat}
 \hat{\mathbf{x}}_r=\mathbf{A}_r^{-1}\big(\mathbf{G}_r^\T\mathbf{C}_d^{-1}\mathbf{d}_r+\mathbf{C}_x^{-1}\bm\mu_r\big).
\end{equation}

For $\beta$ this marginalization is exact. For $t$ and $\theta$ it is a
Laplace (linearized) approximation, which is accurate when $\bm\ell_r$ is
near the conditional mode of $(t,\theta)$ for the current structure.

\subsection{Relinearization}\label{sec:relin}
Every $n_{\rm relin}$ iterations (\texttt{opt.relin}; default 1, i.e.\
before every proposal), each chain state moves its linearization points to
the conditional posterior mean for its current structure. This is one
damped Gauss--Newton step:
\begin{equation}\label{eq:relin}
 \mathbf{b}^0_r\leftarrow\hat{\bm\beta}_r,\qquad
 \mathbf{t}^0_r\leftarrow\mathbf{t}^0_r+\clip(\widehat{\delta\mathbf{t}}_r,\pm\Delta_\tau^{\max}),\qquad
 \bm\theta^0_r\leftarrow\bm\theta^0_r+\clip(\widehat{\delta\bm\theta}_r,\pm\Delta_\theta^{\max}),
\end{equation}
with $\Delta_\tau^{\max}=2$~s and $\Delta_\theta^{\max}=3^\circ$, taken from
Eq.~\eqref{eq:xhat}. Then $\mathbf{G}_r$ and $\ln\mathcal{L}_r$ are
recomputed at the new point.

A structural proposal is always compared with the current state at the
\emph{same} linearization points. Because the linearization points follow
the chain, the sampled distribution is an approximation to the exact
marginal posterior. The approximation error vanishes as the linearization
becomes exact at the conditional mode.

\subsection{Initial linearization points (continuation)}
The initial points $\bm\ell^{(0)}_r$ come from Gauss--Newton iterations of
Eq.~\eqref{eq:relin} at the reference dispersion
$c^w_n=c^{w,\rm ref}_n$, starting from $\mathbf{t}^0=\bm\theta^0=0$ and
$\mathbf{b}^0=1$, with $\sigma_q=\sigma^{\rm true}$ and $r_q=\hat r_q$.
The iterations use a continuation in two stages of 30 steps each:
\begin{enumerate}
\item The higher time-shift terms are held at zero by a tight prior,
  $\mathbf{s}_\tau=(15,10^{-3},10^{-3})$~s, so only a constant $\tau_n$ is
  fitted.
\item All terms are freed with the full prior.
\end{enumerate}
Freeing all terms at once lets overlapping branches trade delays, and
Gauss--Newton then ends in poorer local minima. At the true structure, the
median $\chi^2$ of the Love fundamental windows is 2.0 without the
continuation and 1.13 with it.

The same procedure applied at the true dispersion gives the reference
log-likelihood $\ln\mathcal{L}^{\rm true}$ of the true structure, reported
in the summaries; idagrn data have no ``true'' nuisance parameters.

%======================================================================
\section{Prior on the structure and the noise}\label{sec:prior}
%======================================================================

\begin{itemize}
\item \textbf{Numbers of control points:} uniform on the set
  $\{(k_v,k_x):0\le k_x\le k_v\le k_{\max}\}$, with $k_{\max}=15$
  (with free $\eta$: $\{(k_v,k_x,k_e):0\le k_x,k_e\le k_v\le k_{\max}\}$).
  $\xi$ and $\eta$ never have more interior control points than $\vsv$.
\item \textbf{Control depths:} independent, uniform on $(z_{\min},z_{\max})$.
\item \textbf{Values:} $\vsv$ values (interior and top) uniform on
  $[4.0,5.0]$~km/s; $\xi$ values uniform on $[0.90,1.20]$; $\eta$ values
  (if free) uniform on $[0.80,1.05]$ (\texttt{opt.etamin},
  \texttt{opt.etamax}); bottom values fixed to the reference.
\item \textbf{Noise:} $\ln\sigma_q$ uniform (scale-invariant; implied by
  the symmetric proposal in $\ln\sigma_q$ with acceptance on the
  likelihood ratio alone). $r_q$ uniform on $[0,r_{\max})$, with $r_{\max}=0.98$.
\end{itemize}
The tempered target at temperature $T\ge1$ is
\begin{equation}
 \pi_T(\mathbf{m},\bm\sigma,\mathbf{r}\mid\mathbf{d})\propto p(\mathbf{m})\,p(\bm\sigma)\,p(\mathbf{r})\;\mathcal{L}(\mathbf{m},\bm\sigma,\mathbf{r})^{1/T},
\end{equation}
and $T=1$ is the posterior.

%======================================================================
\section{Sampling}\label{sec:mcmc}
%======================================================================

\subsection{One iteration of one chain state at temperature $T$}
\begin{enumerate}
\item \textbf{Relinearization} (if the iteration number $i$ satisfies
  $(i-1)\bmod n_{\rm relin}=0$): apply Eq.~\eqref{eq:relin} to every record
  and update $\ln\mathcal{L}$.
\item \textbf{Structural move.} With probability $p_\xi=0.4$ the move acts
  on $\xi$, with probability $p_\eta=0.2$ on $\eta$ (only if it is free),
  otherwise on $\vsv$. One of four move types is chosen with
  equal probability:
  \begin{enumerate}
  \item \emph{birth}: add a control point with depth and value drawn from
    the prior (rejected if $k=k_{\max}$, or if it would give $k_x>k_v$ or $k_e>k_v$);
  \item \emph{death}: delete a uniformly chosen interior point (rejected if
    $k=0$, or if it would give $k_v<k_x$ or $k_v<k_e$);
  \item \emph{move}: $z_i\leftarrow z_i+s_z T^{1/4}\varepsilon$, with
    $s_z=10$~km;
  \item \emph{value}: one of the $k$ interior values or the top value,
    chosen uniformly, $v\leftarrow v+s_vT^{1/4}\varepsilon$, with
    $s_v=0.03$~km/s for $\vsv$ and $0.01$ for $\xi$ and $\eta$.
  \end{enumerate}
  Here $\varepsilon\sim\mathcal{N}(0,1)$. A proposal outside the prior
  bounds is rejected. Otherwise the dispersion (Eq.~\ref{eq:disp}), the
  operators $\mathbf{G}_r$ at the current $\bm\ell_r$ and the likelihood
  are recomputed, and the move is accepted with probability
  \begin{equation}
   \alpha=\min\!\Big\{1,\exp\Big[\frac{\ln\mathcal{L}(\mathbf{m}')-\ln\mathcal{L}(\mathbf{m})}{T}\Big]\Big\}.
  \end{equation}
  Birth values come from the prior and deaths are uniform, so the prior and
  proposal terms cancel and only the likelihood ratio remains. The
  $T^{1/4}$ factor widens the steps at high temperature.
\item \textbf{Noise move.} Choose $q$ uniformly. With probability
  $\tfrac12$ propose $\sigma_q'=\sigma_q\,e^{0.05\varepsilon}$, otherwise
  $r_q'=r_q+0.02\varepsilon$. Re-evaluate Eq.~\eqref{eq:logL} for the
  records of that wave type ($\mathbf{G}_r$ is unchanged) and accept with
  the same rule.
\end{enumerate}
Initial states are independent draws from the structural prior, with
$\sigma_q$ equal to the standard deviation of the data of one record of
the wave type, $r_q=0.5$, and $\bm\ell_r=\bm\ell^{(0)}_r$.

\subsection{Plain sampler}
$N_c$ independent chains at $T=1$ (\texttt{opt.sampler='plain'};
$N_c=4$, $N_{\rm it}=3000$ iterations, $1500$ burn-in).

\subsection{Parallel tempering}
$N_{\rm ens}$ independent ensembles (\texttt{opt.sampler='pt'}). Each
ensemble has $K_T$ chains at temperatures
$T_k=T_{\max}^{(k-1)/(K_T-1)}$, $k=1,\dots,K_T$ (log-spaced, $T_1=1$).
After every iteration, swaps of whole states, including their
linearization points, are proposed between neighbouring temperatures
$(p,p+1)$. On alternate iterations these are the pairs with $p$ odd or $p$
even. A swap is accepted with probability
\begin{equation}
 \alpha_{\rm swap}=\min\!\Big\{1,\exp\Big[\big(\ln\mathcal{L}_{p+1}-\ln\mathcal{L}_p\big)\Big(\frac1{T_p}-\frac1{T_{p+1}}\Big)\Big]\Big\}.
\end{equation}
Only the $T=1$ chain is saved. The defaults are $N_{\rm ens}=2$, $K_T=6$,
$T_{\max}=20$ and $1500$ iterations with $750$ burn-in. The current test
uses $K_T=4$, $T_{\max}=10$ and $1200/600$ iterations.

\subsection{Alternative: sampled nuisance parameters}
With \texttt{opt.nuisance='sample'} (plain sampler only), $t_{rna}$ and
$\theta_{rna}$ are Markov-chain variables. The likelihood uses the exact
(non-linearized) waveforms of Eq.~\eqref{eq:forward}, with only $\beta$
marginalized, and their Gaussian prior of Section~\ref{sec:nuisprior}
enters the acceptance ratio. Proposals perturb one coefficient by
$0.4$~s or $0.25^\circ$ times $\varepsilon$.

%======================================================================
\section{Updating the reference model}\label{sec:refupd}
%======================================================================

The reference spectra $U^{\rm ref}_{rjn}$, the reference phase velocities
$c^{w,\rm ref}_n$ and the kernels $G^{w,X}$ of Eq.~\eqref{eq:disp} belong
to one reference model. A reference far from the solution makes the
linearized dispersion inaccurate.

\subsection{At the end of burn-in (default)}
With \texttt{opt.update\_ref = 'burnin'} the run has three phases:
\begin{enumerate}
\item \emph{Burn-in} ($N_{\rm burn}$ iterations) with the step-1 reference.
  The $T=1$ samples of the last fraction $f_{\rm upd}=0.2$
  (\texttt{opt.upd\_frac}) of the burn-in are kept.
\item \emph{Update}, once, after every chain or ensemble has finished its
  burn-in. Let $\bar\vsv(z)$ and $\bar\vsh(z)$ (and $\bar\eta(z)$) be the
  means of these samples over all chains. The new reference card has
  \begin{equation}
   \vsv^{0\prime}=\bar\vsv,\;\; \vsh^{0\prime}=\bar\vsh,\;\;
   \vpv^{0\prime}=R_{pv}\bar\vsv,\;\; \vph^{0\prime}=R_{ph}\bar\vsh,\;\;
   \rho^{0\prime}=R_\rho\bar\vsv\quad\text{on }[z_{\min},z_{\max}],
  \end{equation}
  with $R$ the ratios of the old reference, and $\eta^{0\prime}=\bar\eta$
  if $\eta$ is free. Everything else is unchanged. MINEOS and
  \texttt{idagrn6\_mask} are rerun for it, giving new $U^{\rm ref}$,
  $c^{\rm ref}$ and kernels. The data, windows and normalizations $s_r$
  are kept, so log-likelihoods before and after are evaluated on the same
  data. New initial linearization points $\bm\ell^{(0)}$ are fitted as in
  Section~\ref{sec:relin}, and $\ln\mathcal{L}^{\rm true}$ is recomputed.
\item \emph{Sampling} ($N_{\rm it}-N_{\rm burn}$ iterations). Every chain
  and temperature continues from its own last state $(\mathbf{m},\bm\sigma,\mathbf{r})$.
  Its linearization points are moved to the new reference by
  \texttt{opt.upd\_ngn} $=10$ Gauss--Newton steps from $\bm\ell^{(0)}$.
  The reference is then fixed for the rest of the run.
\end{enumerate}
All chains share one updated reference, so they sample the same target
and $\hat R$ and the parallel-tempering swaps remain valid. Updating
repeatedly \emph{during} sampling would change the target as the chain
moves, so the reference is updated only at the end of burn-in. With
\texttt{opt.update\_ref = 'none'} the step-1 reference is used
throughout.

\subsection{Checking the linearization (D5)}
\texttt{D5\_step2\_check\_linearization.m} evaluates, by full MINEOS and
by Eq.~\eqref{eq:disp} about the run's (post-burn-in) reference:
\begin{itemize}
\item the posterior mean model, and
\item the true model projected on the parameterization (PA5 $\vsv$,
  $\vsh$ on $[z_{\min},z_{\max}]$, Vp and $\rho$ scaled, the reference
  elsewhere).
\end{itemize}
It reports the per-branch rms of four differences:
\begin{itemize}
\item the linearization error at the posterior mean;
\item the linearization error at the projected truth;
\item MINEOS(projected truth) $-$ MINEOS(PA5): the part of the true
  dispersion the parameterization cannot represent, such as PA5 below
  350~km, its Vp/Vs and $\rho$/Vs ratios and $\eta$;
\item MINEOS(posterior mean) $-$ MINEOS(PA5).
\end{itemize}
A reference update can reduce the first two, not the third.

\subsection{Outer iterations (D6, optional)}
\texttt{D6\_step2\_update\_reference\_and\_rerun.m} builds a new reference
from the posterior mean of a finished run, in the same way as the
end-of-burn-in update, and reruns the whole inversion about it. The data
and sampler settings are unchanged, and the tag is \texttt{\_it2}. It can
be repeated until the posterior mean stops changing.

%======================================================================
\section{Outputs and diagnostics}\label{sec:out}
%======================================================================

After burn-in every 5th iteration is saved (thinning). Each sample stores:
\begin{itemize}
\item $\vsv$, $\vsh$, $\xi$ and $\eta$ on a 2-km depth grid (0--400~km);
\item the dispersion curves $c^w_n(T_\ell)$;
\item $k_v$, $k_x$, $\bm\sigma$ and $\mathbf{r}$;
\item the conditional posterior means of the nuisance coefficients,
  $\mathbf{t}^0+\widehat{\delta\mathbf{t}}$ and
  $\bm\theta^0+\widehat{\delta\bm\theta}$, per record and branch.
\end{itemize}
Convergence is measured with the Gelman--Rubin potential scale reduction
factor of each scalar quantity (e.g.\ $\vsv$ at each depth) across chains.
For $M$ chains of $n$ samples, with within-chain variance $W$ (the mean of
the chain variances) and between-chain variance $B=n\,\mathrm{Var}(\text{chain means})$,
\begin{equation}
 \hat R=\sqrt{\frac{\frac{n-1}{n}W+\frac{B}{n}}{W}},
\end{equation}
and $\hat R\to1$ at convergence.

The waveform figures show, for each record:
\begin{itemize}
\item the data $\mathbf{d}_r$;
\item the PREM synthetic;
\item the reference prediction $\Ifft[\sum_nU^{\rm ref}_{rjn}]$, with no
  nuisance correction;
\item the final prediction: the exact waveforms of Eq.~\eqref{eq:forward}
  for the last chain state, with $t$ and $\theta$ set to their means over
  the last 20 samples and $\beta$ to its conditional posterior mean.
\end{itemize}

%======================================================================
\section{Approximations and remarks}
%======================================================================
\begin{itemize}
\item The dispersion is linearized about a reference model
  (Eq.~\ref{eq:disp}). MINEOS and \texttt{idagrn6} are rerun only when the
  reference is updated (Section~\ref{sec:refupd}), never for individual
  proposals.
\item The array is treated as flat and the wavefield as locally plane
  across it. Amplitude variations across the array come only from
  $U^{\rm ref}$.
\item Data branches $n_m,\dots,N_b^{\rm data}-1$ (8--15) are not modeled.
  Their energy in the overtone windows acts as additional noise, which the
  hierarchical $\sigma_q$ and $r_q$ absorb.
\item The time shifts $\tau_n(\omega)$ and arrival angles
  $\delta\theta_n(\omega)$ are common to all stations of a record. Their
  variation across the array is described by the structure only.
\end{itemize}

%======================================================================
\section{Notation}\label{sec:notation}
%======================================================================
\begin{longtable}{>{$}l<{$}p{0.68\textwidth}}
\toprule
\text{Symbol} & Meaning \\
\midrule
\endhead
e,\;N_e & event index; number of events (15)\\
j,\;N_s & station index; number of stations (16)\\
w & wave type: $L$ (Love, transverse) or $R$ (Rayleigh, vertical)\\
r=(e,w),\;N_r & record (one event, one component, all stations); number of records (30)\\
n,\;n_m & mode branch index ($0$ = fundamental); number of modeled branches (8)\\
N_b^{\rm data} & number of branches in the synthetic data (16)\\
q & window type: 1 Love overtone, 2 Love fundamental, 3 Rayleigh overtone, 4 Rayleigh fundamental\\
(x_j,y_j) & station east/north coordinates relative to the array center (km)\\
\Delta_e,\;\Delta_{ej} & epicentral distance to the array center and to station $j$ (km)\\
\varphi_e & propagation azimuth at the array (deg clockwise from north)\\
\Delta D_j(\varphi) & plane-wave path difference of station $j$ relative to the array center (km)\\
\Delta t,\;N,\;t_i,\;t^0_r & sampling interval (4 s), segment length (800), sample times, segment start\\
f_k,\;\omega_k,\;\mathcal{B},\;N_f & frequencies, angular frequencies, band, number of band frequencies\\
f_1\dots f_4 & band-pass corners (1/112, 1/100, 1/20, 1/17 Hz)\\
h,\;H & time-domain Tukey window; frequency-domain cosine band-pass taper\\
\Ifft & band spectrum $\to$ real time series, Eq.~\eqref{eq:ifft}\\
U^{w,\rm ref}_n(T) & reference-model group velocity of branch $n$\\
U^{\rm b}_w,U^{\rm ov}_w,U^{\rm f}_w & group-velocity bounds of the windows (km/s)\\
\mathbf{d}_r,\;N_{d,r},\;N_d & data vector of record $r$; its length; total number of data\\
s_r & normalization of record $r$ (peak of the noise-free windowed data)\\
\sigma^{\rm true},\;\hat r_q & noise level of the synthetic noise (0.03); empirical lag-one correlation\\
U^{\rm ref}_{rjn}(\omega) & reference spectrum: single-branch synthetic of the reference model\\
z,\;z_w & depth below the seafloor; water depth (km)\\
z_{\min},z_{\max} & inverted depth range (7, 350 km)\\
\vsv,\vsh,\vpv,\vph,\rho,\eta & TI elastic parameters\\
\xi & radial anisotropy $(\vsh/\vsv)^2$\\
\mathbf{m} & structural model: control points of the $\vsv$ and $\xi$ B-splines\\
k_v,k_x,k_{\max} & numbers of interior control points of $\vsv$ and $\xi$; maximum (15)\\
(z^p_i,v^p_i) & interior control points of profile $p$\\
v^p_{\rm top},v^p_{\rm bot} & end values at $z_{\min}$ (free) and $z_{\max}$ (fixed)\\
N_{l,\varrho}(s),\;\varrho & B-spline basis functions; degree ($\le3$)\\
R_{pv},R_{ph},R_\rho & reference-model ratios $\vpv/\vsv$, $\vph/\vsh$, $\rho/\vsv$\\
\eta,\;k_e,\;(z^e_i,v^e_i) & anisotropic parameter $F/(A-2L)$; its interior control points (if free)\\
\mathcal{X},\;\delta X,\;X^0 & the six TI parameters; perturbation of $X$ from the reference; reference value\\
T_\ell,\;N_T & dispersion-node periods; their number (10)\\
c^w_n(T;\mathbf{m}),\;c^{w,\rm ref}_n & model and reference phase velocities (km/s)\\
\mathcal{N} & card nodes in the inverted depth range\\
K^{w,n}_X(T,r) & MINEOS Fréchet kernel of phase velocity w.r.t.\ parameter $X$\\
\Delta r_i & node weight (m)\\
G^{w,X} & kernel of parameter $X$ at the dispersion nodes and model nodes\\
\chi^w_n(\omega) & mask: 1 where branch $n$ exists in the band\\
B_{rn}(\omega),\;\beta_{rnk},\;\Phi_k,\;K & spectral correction; its coefficients; hat functions in $\ln\omega$; their number (3)\\
\tau_{rn}(\omega),\;t_{rna},\;K_\tau & time shift (s); Legendre coefficients; number of terms (3)\\
\delta\theta_{rn}(\omega),\;\theta_{rna},\;K_\theta & arrival-angle anomaly (deg); coefficients; number of terms (2)\\
\zeta(\omega),\;P_a & normalized log frequency on $[-1,1]$; Legendre polynomials\\
\mathbf{x}_r,\;\tilde{\mathbf{x}}_r,\;P & nuisance vector; its linearized form; its length (88)\\
\bm\ell_r=(\mathbf{b}^0,\mathbf{t}^0,\bm\theta^0) & linearization point of record $r$\\
S_{rjn}(\omega) & branch spectrum at the linearization point\\
\mathbf{G}_r,\;\mathbf{G}^\beta_r & linearized operator; its amplitude columns\\
\tilde{\bm\beta}_r,\;\lambda,\;\bar\sigma & damped least-squares amplitudes used in the derivative columns; damping; mean noise level\\
\bm\mu_r,\;\mathbf{C}_x & prior mean and covariance of $\tilde{\mathbf{x}}_r$\\
\sigma_b,\;\mathbf{s}_\tau,\;\mathbf{s}_\theta & prior std of $\beta$ (0.4), of $t$ ((15,10,10) s), of $\theta$ ((5,5)$^\circ$)\\
\sigma_q,\;r_q,\;r_{\max} & noise level and lag-one correlation of window type $q$; upper bound of $r_q$ (0.98)\\
\mathbf{C}_{d,r} & AR(1) noise covariance of record $r$\\
n_{rq},\;m_{rq} & number of samples, number of contiguous segments of type $q$ in record $r$\\
\mathbf{A}_r,\bm\rho_r,\mathbf{g}_r & posterior precision of $\tilde{\mathbf{x}}_r$; prior residual; projected residual\\
\mathcal{L}_r,\;\mathcal{L} & marginal likelihood of record $r$; total\\
\hat{\mathbf{x}}_r & conditional posterior mean of $\tilde{\mathbf{x}}_r$\\
n_{\rm relin} & iterations between relinearizations (default 1)\\
\Delta_\tau^{\max},\Delta_\theta^{\max} & maximum relinearization steps (2 s, 3$^\circ$)\\
T,\;T_k,\;T_{\max},\;K_T & temperature; ladder; maximum; number of temperatures\\
\pi_T & tempered posterior\\
p_\xi,\;p_\eta,\;s_z,\;s_v & probabilities of a $\xi$ (0.4) and an $\eta$ move (0.2, if free); depth step (10 km); value steps\\
N_{\rm burn},\;f_{\rm upd} & burn-in iterations; fraction of the burn-in averaged for the reference update (0.2)\\
\varepsilon & standard normal random number\\
\alpha,\;\alpha_{\rm swap} & Metropolis acceptance probabilities\\
N_c,\;N_{\rm ens},\;N_{\rm it} & number of chains, ensembles, iterations\\
\hat R,\;W,\;B & Gelman--Rubin factor; within- and between-chain variances\\
\bottomrule
\end{longtable}

\end{document}
